\section{Trained and semi-random networks}\label{sec:trained}

The theorems so far assume Gaussian weights. Trained networks are not Gaussian, but empirically their
weight matrices are a Gaussian bulk plus a small number of outlier directions, and the balance route to
the small output does not need Gaussianity at all. This section records what survives training: a
perturbation-stability statement that degrades the Gaussian bounds gracefully under a bounded deviation
of the weights, the balance version of the fragility theorem for trained classifiers, and a discussion
of the numerical evidence.

\subsection{Perturbation stability}
\begin{theorem}[Stability under weight deviation]\label{thm:trained}
Let $W_\ell=W_\ell^0+E_\ell$ with $W_\ell^0$ Gaussian (satisfying the hypotheses of \Cref{thm:master})
and $E_\ell$ arbitrary with $\|E_\ell\|_\op\le\eps_\ell$. Then:
\begin{enumerate}
\item the Lipschitz certificate of \Cref{thm:certificate} holds with $L_{\mathrm{net}}$ replaced by
$\|w\|\prod_\ell(2\sqrt{\kappa(c_1)}+\eps_\ell)$;
\item the gradient lower bound of \Cref{prop:gradlb} survives, $\|\nabla F(x)\|\ge c_g\|w\|-\sum_\ell\eps_\ell G_\ell$,
where $G_\ell$ is the product of the operator-norm bounds of the other layers, provided the deviation
is small, $\sum_\ell\eps_\ell G_\ell<\tfrac12 c_g\|w\|$;
\item consequently the fragility conclusion of \Cref{thm:master} holds for the perturbed network under
this smallness condition, with the constants shifted accordingly.
\end{enumerate}
\end{theorem}

\begin{proof}
(1) Operator norms are subadditive, $\|D_\ell(W_\ell^0+E_\ell)P\|_\op\le\|D_\ell W_\ell^0P\|_\op+\eps_\ell$,
so each factor's bound increases by $\eps_\ell$ and the product bound is
$\prod_\ell(2\sqrt{\kappa(c_1)}+\eps_\ell)$. (2) Write the Jacobian of the perturbed network as the
Gaussian Jacobian plus a sum of $L$ terms, each replacing one $W_\ell^0$ by $E_\ell$ and bounding the
other factors by their operator norms; the correction is at most $\sum_\ell\eps_\ell G_\ell$ in norm,
which is subtracted from the Gaussian lower bound $c_g\|w\|$. (3) Under the smallness condition the
perturbed gradient is still $\ge\tfrac12c_g\|w\|$, and the small output and stability are perturbed by
the same $O(\sum\eps_\ell)$ amounts, so \Cref{prop:engine} applies with the shifted constants.
\end{proof}

The smallness condition is a low-rank/small-norm deviation from the Gaussian bulk, which is what the
Marchenko--Pastur analysis of trained weights \cite{shmaloRMT,shmaloAccuracy} identifies: the bulk of
the singular spectrum stays at the random edge while a few directions carry the learned signal, so
$\|E_\ell\|_\op$ is the top learned singular value, typically $O(1)$ and not growing with width, while
$c_g\|w\|$ and the certificate scale are $\Theta(1)$. The theorem then says training moves the
constants but not the $1/\sqrt n$ scale.

\subsection{The balance route for trained classifiers}
For a trained classifier the readout is not random, so the small output must come from balance. If the
trained classifier is accurate and its ground-truth labels are balanced, then its sign is balanced under
the input measure (pushed to the sphere), and \Cref{prop:balance} gives the small output for all but a
$\delta$-fraction of inputs. The gradient facts, however, need the weights to be near-Gaussian, which is
\Cref{thm:trained}. Combining:

\begin{corollary}[Trained-network fragility, informal]\label{cor:trained-frag}
Let $F$ be a trained classifier on the sphere whose weights are a Gaussian bulk plus a bounded deviation
(\Cref{thm:trained}) and whose sign is $p$-balanced. Then for all but a $\delta$-fraction of inputs a
single gradient step of length $O(\eps_{\mathrm{adv}}(p,\delta)\cdot L_{\mathrm{net}}/c_g)$ crosses the
boundary.
\end{corollary}

This is the trained-regime statement of \Cref{thm:master}: the small output is measure-theoretic (from
balance) rather than pointwise (from a random readout), and the gradient facts hold up to the weight
deviation. It is the theorem that the CIFAR experiments of \cite{advExt} were designed to probe.

\subsection{Numerical evidence from trained networks}
The experiments of \cite{advExt} train small CNNs, ResNets, and ViTs on CIFAR-10 and CIFAR-100 and
measure, along training, the one-step attack success, the predicted radius $2m/\|\nabla m\|$ against the
empirical attack radius, and the compatibility ratio. Three of their findings are exactly what the
present theory predicts and are worth restating in its language. First, the empirical attack radius is
monotone in the predicted radius $2m/\|\nabla m\|$ and stays below it, on both datasets --- the
conservative direction, since the prediction is an upper bound from the small-output estimate. Second,
the one-step vulnerability signatures appear early and persist after the networks become accurate, so
fragility is not a symptom of poor accuracy but a structural property of the map, as the balance-plus-%
gradient mechanism requires. Third, the transformer shows the strongest signatures despite the lowest
accuracy, consistent with the near-identity residual structure making its gradient lower bound most
robust. The experiments are diagnostics, not certified radii, and the frozen artifacts are catalogued in
\cite{advExt}; we use them as the trained-regime evidence for the scale that \Cref{thm:master} and
\Cref{cor:sharp} predict, and add the new numerical verification of the product rate in the next section.
